\selectlanguage{french}
\chapter*{Résumé}
\addcontentsline{toc}{chapter}{Résumé}

\noindent Pendant 38 mois d'alternance chez SAS Everbloo (octobre 2023 -- novembre 2026), j'ai travaillé comme apprenti architecte logiciel sur la plateforme \textbf{Metis}. Metis est un moteur de réservation de billets d'avion B2B utilisé par des réseaux d'agences comme TourCom et des entreprises pour leurs déplacements professionnels. J'ai contribué aux quatre dépôts du projet (\texttt{api-aerial}, \texttt{front-reservation}, \texttt{front-dashboard} et \texttt{api-dashboard}), avec 1\,554 commits sur l'ensemble de la stack. Le principal défi technique a été de faire fonctionner ensemble deux mondes très différents : les GDS historiques (Sabre, Amadeus), qui reposent sur de vieilles sessions avec état et le protocole EDIFACT, et les API modernes IATA NDC, qui fonctionnent en requêtes web sans état avec des offres valables quelques minutes. C'est tout le sujet de ce mémoire : comment garantir que les prix et les réservations restent cohérents quand on mélange des technologies aussi différentes ?

\noindent Pour résoudre ce problème sur le terrain, j'ai mis en place une architecture modulaire avec un backend Node.js/Express (ORM Sequelize) et des interfaces Nuxt~3 / Vue~3 en TypeScript. J'ai d'abord créé une couche d'adaptateurs (\texttt{SabreAdapter.ts}) pour convertir toutes les réponses des compagnies vers un format unique, et un intercepteur Axios pour renouveler les sessions OAuth2 en arrière-plan sans déconnecter les utilisateurs. J'ai aussi développé un algorithme (\texttt{matchesPaxRefId}) pour régler les rejets fréquents de réservations IATA, notamment sur les bébés sans siège ou les noms avec accents. Quand les prix changent au dernier moment, j'ai mis en place un dialogue de confirmation interactif au lieu de bloquer l'achat, et j'ai corrigé la gestion multi-devises (norme ISO 4217) pour supprimer les erreurs de calcul de centimes et les pénalités financières (ADM).

\noindent Ces solutions ont été éprouvées sur le trafic réel de près de cent agences de voyages. Sur la période observée en production, nos logs montrent un taux de rejet applicatif NDC inférieur à 0,1\,\% (hors indisponibilités côté transporteurs), le temps de réponse moyen des recherches est tombé à 3,0\,s lors de nos campagnes de tests, et nous n'avons eu aucun litige financier (ADM) sur les deux dernières années. C'est surtout en production que j'ai compris que ces problèmes théoriques avaient des conséquences très concrètes : décalages de fuseaux horaires, deadlocks SQL lors de pics de trafic ou gestion du stress lors des pannes du lundi matin. Cette alternance m'a permis de passer du rôle de développeur qui résout des tickets à celui d'ingénieur capable de concevoir et maintenir une plateforme critique.

\vspace{0.6cm}
\noindent \textbf{\large Mots-clés :}\\
\noindent Distribution Aérienne, GDS (Sabre, Amadeus), IATA NDC, Nuxt 3, Vue 3, TypeScript, Node.js, Express, Cohérence Transactionnelle, Reshopping, OAuth2, Multi-Devises ISO 4217, Segments Passifs, Tests Automatisés.
